% Compiler: pdfLaTeX
\documentclass[11pt,a4paper]{article}

\usepackage[utf8]{inputenc}
\usepackage[T1]{fontenc}
\usepackage{lmodern}
\usepackage{microtype}
\usepackage{geometry}
\usepackage{enumitem}
\usepackage{titlesec}
\usepackage{fancyhdr}
\usepackage{hyperref}

% Document details
\newcommand{\researchtitle}{Path dependence}
\newcommand{\researchshorttitle}{Path dependence}
\newcommand{\researchdate}{2026-10-25}

% Layout
\geometry{left=2.7cm,right=2.7cm,top=2.5cm,bottom=2.5cm,headheight=14pt}
\hypersetup{hidelinks,pdftitle={\researchtitle}}
\urlstyle{same}
\frenchspacing
\setlength{\parindent}{0pt}
\setlength{\parskip}{0.55em}
\linespread{1.04}
\setlength{\emergencystretch}{1em}

\titleformat{\section}
    {\large\bfseries}{\thesection}{0.65em}{}
\titleformat{\subsection}
    {\normalsize\bfseries}{\thesubsection}{0.65em}{}
\titlespacing*{\section}{0pt}{1.15em}{0.35em}
\titlespacing*{\subsection}{0pt}{0.7em}{0.2em}
\setlist[itemize]{leftmargin=1.35em,itemsep=0.25em,topsep=0.2em}
\setlist[enumerate]{leftmargin=1.6em,itemsep=0.25em,topsep=0.2em}

\pagestyle{fancy}
\fancyhf{}
\fancyhead[L]{\small Random Research}
\fancyhead[R]{\small\researchshorttitle}
\fancyfoot[C]{\small\thepage}
\renewcommand{\headrulewidth}{0pt}
\renewcommand{\footrulewidth}{0pt}

\begin{document}

{\LARGE\bfseries\researchtitle\par}
\vspace{0.4em}
{\small\researchdate\par}

\section{Overview}

Path dependence describes processes in which earlier events help shape
the outcomes available later. Once people have built skills and routines
around one option, a different starting point may no longer lead to the
same result. \cite{arthur,david}

The idea makes the history of a choice relevant. Knowing what alternatives
exist today may not explain how one became established.

\section{Findings}

Arthur models competing technologies that improve as more people adopt
them. In that setting, small historical events can influence which
technology becomes established. The outcome can become difficult to
change and need not be the most efficient one. This is a result of the
model's assumptions about increasing returns, not a claim about every
technology market. \cite{arthur}

David's account adds a useful limit. Path dependence alone is neither
necessary nor sufficient for market failure. Historical influence does
not establish that the outcome is inferior, or that a particular
intervention would improve it. \cite{david}

Together, these accounts make a narrow argument: the route to an outcome
can matter, while its quality still needs a separate assessment.

\section{Example}

Imagine a small team chooses a spreadsheet tool because one early member
already knows it. Over time, the team builds templates, trains new members,
and exchanges files with partners using the same format.

A competing tool later offers features they prefer. Switching would still
require converting files and learning new workflows. The original
decision has changed the cost of choosing today.

This hypothetical example also helps separate path dependence from the
sunk cost fallacy. Time already spent learning is gone. The work that a
migration would require is still a future cost and belongs in the decision.

\section{Discussion}

I find this more convincing when somebody can explain what keeps the
current arrangement in place. Saying that history matters is too broad
on its own. In the example, the templates and shared skills provide
specific reasons for persistence.

I would also be careful about assuming that an old arrangement survives
only because people resist change. Compatibility and familiarity may
still be useful. A newer option has to be assessed alongside the effort
required to move.

For me, the idea raises a question about early choices: which small
conveniences could become expensive commitments if the project grows?

\section{Practical implications}

For a tool expected to last, I would check how easily its files can be
exported and whether a small trial could reveal problems before others
depend on it.

When reviewing an established choice, I would separate the reason it
was first adopted from the reasons to keep it now. Some may have expired.
Others may remain strong enough to justify staying.

\renewcommand{\refname}{References}
\begin{thebibliography}{9}
\small
\setlength{\itemsep}{0.35em}

\bibitem{arthur}
Arthur, W. B. (1989).
\href{https://doi.org/10.2307/2234208}
{Competing technologies, increasing returns, and lock-in by historical events.}
\textit{The Economic Journal}, 99(394), 116--131.

\bibitem{david}
David, P. A.
\href{https://www.nuffield.ox.ac.uk/economics/history/paper20/path.htm}
{Path dependence and the quest for historical economics:
One more chorus of the ballad of QWERTY.}
Author's abstract, University of Oxford.

\end{thebibliography}
\end{document}
